% Chronological predecessor studies; statistics are version-scoped.
\section{V2：以过程一致性为中心的可追溯工作流}
\label{sec:v2}
\subsection{研究思路}
V2 首先处理实际科学计算中过程与证据的一致性问题。其基本组织方式是将已知计算协议分解为明确步骤，在受控启动、状态保存和产物关联的基础上形成可复核计算链。模型决策并非这一阶段量化成果的主要来源；研究重点在于确认计算确实发生、原始结果得到保留，并能够在不重算的情况下恢复其科学含义。

这一设计对应自主计算的基础要求：若输入几何、实际执行条件或输出来源不能被可靠识别，即使后续加入语言模型，也难以判断模型所使用的数值是否属于正确任务。因此，V2 将真实执行与离线科学重解析分离，以固定结果验证证据保存和解释的一致性。

\subsection{计算结果}
本文采用两条人工逐节点批准的成功计算链作为 V2 案例。水和乙醇均使用 ORCA 6.1.1、气相、零电荷、单重态及 r$^{2}$SCAN-3c/TightSCF，执行优化、频率和单点计算。r$^{2}$SCAN-3c 是具有组合校正的复合电子结构方法\cite{r2scan3c}。后续科学分析读取已有输出，没有启动重复计算。结果见表 \ref{tab:v2} 和图 \ref{fig:v2-convergence}。

\begin{table}[htbp]
\centering\small
\caption{V2 两条受控计算链的结果。耗时为 ORCA 各节点自报经过时间之和，未包含人工等待及工作流开销；不同分子和核数不用于计算加速比。}
\label{tab:v2}
\begin{tabularx}{\linewidth}{p{5.3cm}YY}
\toprule
指标 & 水 H$_2$O & 乙醇 C$_2$H$_6$O\\
\midrule
核数 / 总内存 / MaxCore & 1 / 2048 MB / 1536 MB & 4 / 2048 MB / 384 MB\\
成功计算节点 & Opt、Freq、SP 各一次 & Opt、Freq、SP 各一次\\
优化周期 & 4 & 11\\
优化末次评价电子能 / $\Eh$ & $-76.418938721015$ & $-155.002314096955$\\
后续单点电子能 / $\Eh$ & $-76.418938721035$ & $-155.002314094574$\\
首周期至优化终态降能 / m$\Eh$ & 1.692636839 & 7.859199377\\
正振动模式数 & 3 & 21\\
正频率范围 / cm$^{-1}$ & 1653.249--3932.735 & 248.134--3837.522\\
三节点自报总耗时 / s & 24.143 & 98.855\\
\bottomrule
\end{tabularx}
\end{table}

\begin{figure}[htbp]
\centering
\includegraphics[width=\linewidth]{v2-optimization.pdf}
\caption{V2 水与乙醇的优化收敛曲线。纵轴为各周期电子能减去对应优化末次评价能量，单位 m$\Eh$。收敛后的末次评价仅提供零点，不额外计作优化周期。}
\label{fig:v2-convergence}
\end{figure}

水与乙醇分别获得 3、21 个正振动模式，历史检查支持其给定理论水平下的局部极小值。定义 $\Delta E=E_{\mathrm{SP}}-E_{\mathrm{Opt,final}}$，两者的内部重复评价差分别为 $-2.0\times10^{-11}\,\Eh$ 与 $2.381\times10^{-9}\,\Eh$。该差值反映同一终态的数值一致性，而非相对实验或高精度参考的误差。

\subsection{方法边界与后续问题}
上述成果建立了实际计算与证据复核的可行性，但不构成无人干预的成功率估计：两条链经过人工批准，水链的结果采集还包含开发者修正。单一初始构型也不能支持全局构象结论。V2 因此提供的是可靠操作与可复核证据的基础，而尚未检验自然语言需求如何自主转化为组合计算。后者构成 V3 的主要研究问题。

\section{V3：有限任务范围内的语义驱动工具组合}
\label{sec:v3}
\subsection{研究思路}
V3 将研究对象从固定计算链扩展到自然语言交互与工具组合。系统在已定义的科学能力范围内解释需求，建立结构化请求，并通过组合工具完成单点、优化、频率和派生量分析。相较 V2，其重点从过程保存转向用户语义与可执行操作之间的连接。

这一思路将科学求解能力主要保留在既有工具中，模型参与语义规范化、结果回答和有限修复。因而，V3 的评价既需验证基础工具的科学输出，也需考察对话情境中的条件修改、已有结果查询、澄清、拒绝与失败恢复。两类评价应区分统计。

\subsection{固定科学矩阵}
Scientific Matrix v1 在 5 类任务与 6 个分子对象中显式选择 11 个单元，11/11 单元通过，实际 ORCA 为 16 次、模型调用为 0，整体经过时间为 421.516 s。该矩阵采用固定请求和计划，检验组合执行与科学检查，不作为模型驱动端到端能力的证据。覆盖范围和部分科学结果见图 \ref{fig:v3-matrix}、表 \ref{tab:v3-science}。

\begin{figure}[htbp]
\centering
\includegraphics[width=\linewidth]{v3-scientific-matrix.pdf}
\caption{V3 固定科学矩阵。11 个启用单元全部通过，灰格未纳入评价分母。单元内标注经过时间；该批次使用真实 ORCA，模型调用为 0。}
\label{fig:v3-matrix}
\end{figure}

\begin{table}[htbp]
\centering\small
\caption{V3 科学矩阵的优化与频率结果。理论水平为气相 r$^{2}$SCAN-3c、零电荷、单重态，资源为 4 核、1024 MB、MaxCore 192 MB/进程。}
\label{tab:v3-science}
\begin{tabularx}{\linewidth}{p{1.8cm}rp{4.6cm}Y}
\toprule
分子 & 优化周期 & 最终优化电子能 / $\Eh$ & 相应频率检查\\
\midrule
H$_2$O & 5 & $-76.418938720985$ & 3 个正模式，局部极小值支持\\
NH$_3$ & 5 & $-56.541874067217$ & 6 个正模式，局部极小值支持\\
CO$_2$ & 3 & $-188.570062871358$ & 4 个正模式，局部极小值支持\\
CH$_4$ & 2 & $-40.494915367074$ & 仅优化，未作频率结论\\
乙醇 & 36 & $-155.002354316872$ & 仅优化，未作频率结论\\
\bottomrule
\end{tabularx}
\end{table}

水的三个正振动频率为 1653.24、3813.59、3932.75 cm$^{-1}$；CO$_2$ 为 668.14、668.14、1362.20、2422.54 cm$^{-1}$。线性 CO$_2$ 的五个外部零模式与四个振动模式被分别检查。两次固定矩阵记录的对照中，13 个电子能结果在输出的 12 位小数精度上相同，三套频率一致。这支持指定输入和设置下的数值重复性，不替代外部准确度评价。

\subsection{模型参与的核心基准}
Benchmark 1.2 core 的最终评分为 18/18，其分母包括 11 个固定科学题、5 个脚本化对话情境和 2 个名称解析题。16 项实际涉及计算，另两项为澄清与拒绝。实验驱动发送预定义话语，并在公开确认状态发送确认回复，经系统的正常对话入口执行，没有直接注入规范请求或计划。

该批次使用 DeepSeek Flash，共 37 次模型调用：语义规范化 19 次、结果回答 17 次、修复 1 次；专门 Intake/Planner 模型调用均为 0。输入、输出和总 tokens 分别为 140,953、6,985 和 147,938，整体经过时间为 667.547 s。ORCA 共尝试 23 次，包括 22 次成功和 1 次预设受控失败，随后执行一次有界修复。图 \ref{fig:v3-cost} 给出调用与耗时分布。

\begin{figure}[htbp]
\centering
\includegraphics[width=\linewidth]{v3-core-costs.pdf}
\caption{V3 核心基准的模型调用分工与任务组耗时。18/18 表示预定义情境的最终评分，其中两项不执行计算；23 次 ORCA 尝试包含一次预设失败及其后续修复。各组的对象和步骤数不同，耗时不作效率排序。}
\label{fig:v3-cost}
\end{figure}

受控失败通过限制优化迭代次数产生，用于检验已定义修复路径，不能解释为自然故障发生率。最终评分还对一项确定性拒绝情境作过规则校准，保存观察没有重跑。因此，18/18 应理解为有限范围内的情境评分，而非任意自然语言任务的自主成功概率。

\subsection{方法边界与 V4 的研究动机}
V3 说明，自然语言语义、工具执行和有界修复可以在有限任务集合中协同运行。然而，科学题的方法、结构和目标受到预定义条件约束，对话也由脚本组织；完整扩展集与未见任务集未运行，数值结果未相对外部高精度参照评价。这些限制使现有成果主要支持已定义范围内的组合能力。

在此基础上，新的问题是：当科学需求要求已有工具未封装的采样、分析或算法时，是否必须由开发者增加专用求解模块？本文据此将 V4 的研究重点转向任务算法的现场生成，通过固定通用工具与生成代码的组合拓展任务表达，同时保持执行边界和科学验证要求。该过渡是基于已有证据的研究问题重构，不构成版本性能提升的预先结论。
